\documentclass{situationist-preprint}

\title{No Permanent Self}
\msid{PID-13-a}
\msfamily{PID-13}
\msdiscipline{Personal Identity and the Self}
\mstype{Principal manuscript}
\mskeywords{self; no-self; process; Hume; minimal self; animalism}

\begin{document}
\maketitle

\begin{abstract}
Is there a permanent self: something that is the subject of a life and remains the same while all its experiences, thoughts and actions come and go? We argue that there is not. The self is a process, the ongoing integration of experience, memory and agency around a first-person perspective, and it persists only as long as the activities that constitute it. Whatever plays the role of the self is either a persisting bearer of the integrating activities or the activity itself. No bearer distinct from the organism is found in experience or needed to explain the unity the activity already explains; and the organism, which on PID-05 thinks your thoughts, exists before the activity begins and can outlive it, so it is not a permanent self. This is consistent with PID-05: you are an organism, and your self is a process you undergo, as your life is. It is also consistent with PID-04, since a process is as real as a thing.
\end{abstract}

\section{The problem}

Several traditions affirm a permanent self: the Upanishadic \emph{\=atman}, the Cartesian thinking substance, the simple soul of rational psychology. Two traditions deny it. Buddhist philosophers argue that what we call a self is a series of impermanent psycho-physical elements to which the label is conventionally applied \citep{Siderits2003}; the companion manuscript PID-13-b relates that view to Parfit's. \citet[1.4.6]{Hume1739} reported that when he entered most intimately into what he called himself he always stumbled on some particular perception and never caught himself without one, and concluded that the self is a bundle of perceptions. Recent views make the self a construct: a self-model \citep{Metzinger2003} or a center of narrative gravity \citep{Dennett1992}. \citet{Strawson2009} argues that selves exist but are short-lived, each lasting no longer than an unbroken stretch of experience. \citet{Zahavi2005} defends a minimal self, the first-personal character of every experience, which he holds to be real and not a construct.

We argue that there is no permanent self and that the self is a process. Since PID-05 holds that you are a thinking organism, something does persist through your experiences, namely you; we show that this thinker is not a permanent self.

\section{Framework}\label{s:frame}

\begin{definition}\label{d:role}
The \emph{self-role} is the role of that which unifies a subject's simultaneous experiences into one point of view, links her experiences, memories and intentions across time as hers, and is what she refers to in first-person thought as the one who perceives, remembers and acts.
\end{definition}

\begin{definition}\label{d:perm}
A \emph{permanent self} is something that plays the self-role and does so for as long as it exists, remaining the same through every change in, and every cessation of, the activities of perceiving, remembering and acting.
\end{definition}

\begin{definition}\label{d:process}
A \emph{process} is something that consists of activities unfolding over time and exists at a time only if some of its constituent activities are then going on.
\end{definition}

\begin{assumption}\label{a:explain}
The unity of simultaneous experiences, the linking of experiences, memories and intentions across time, and first-person reference are explained by integrating activities of the subject: the binding of perceptual information, the encoding and retrieval of memory, the formation and execution of intentions, and the maintenance of a self-representation by which the subject represents these states as its own.
\end{assumption}

\begin{assumption}\label{a:hume}
No experience presents a bearer of experiences that is distinct from the organism and persists through every change in its experiences.
\end{assumption}

\begin{assumption}\label{a:org}
The thinker of a human being's thoughts is a human organism (PID-05), which exists before its integrating activities begin and can survive their irreversible end.
\end{assumption}

Assumption~\ref{a:explain} summarizes the explanatory practice of cognitive science; it does not assume that the self is an illusion. Assumption~\ref{a:hume} is Hume's observation. Assumption~\ref{a:org} combines the conclusion of PID-05-a with the facts of development and of the persistent vegetative state, in which the organism lives on after the integrating activities have permanently ceased.

\section{The argument}\label{s:arg}

\begin{lemma}\label{l:dich}
Whatever plays the self-role is either a bearer of the integrating activities, which persists through their changes, or the integrating process itself.
\end{lemma}

\begin{proof}
By Assumption~\ref{a:explain}, the self-role is filled by the integrating activities. Something plays the role either because it is the subject of those activities or because it consists of them; in the first case it is a bearer, and in the second, by Definition~\ref{d:process}, a process.
\end{proof}

\begin{lemma}\label{l:nobearer}
No bearer other than the organism plays the self-role.
\end{lemma}

\begin{proof}
By Assumption~\ref{a:hume}, experience does not present such a bearer, and by Assumption~\ref{a:explain}, what the self-role requires is explained without it. Moreover, by Assumption~\ref{a:org} the thoughts already have a thinker, the organism, so that a further bearer would be a second thinker of the same thoughts, which PID-05-a rejects.
\end{proof}

\begin{lemma}\label{l:organism}
The organism is not a permanent self.
\end{lemma}

\begin{proof}
By Assumption~\ref{a:org}, the organism exists before its integrating activities begin and can survive their irreversible end. During those periods nothing unifies its experiences or links its memories, since it has none; so it does not then play the self-role. It does not play the role for as long as it exists, and by Definition~\ref{d:perm} it is not a permanent self.
\end{proof}

\begin{theorem}\label{t:main}
There is no permanent self. The self is the integrating process, and it exists only as long as its constituent activities.
\end{theorem}

\begin{proof}
By Lemmas~\ref{l:dich}--\ref{l:organism}, the only candidate bearer, the organism, is not a permanent self, and no other bearer plays the self-role. So what plays the role is the integrating process, and by Definition~\ref{d:process} it exists only while some of its constituent activities go on. Such a process does not remain through their cessation, so it is not a permanent self either.
\end{proof}

You are the organism; your self is a process the organism undergoes, as your life is, and you could outlive it. Your self is nonetheless real: PID-04 holds that the self is a construction and that constructions are real, and the integrating process is such a construction, built and maintained by the organism's own activity.

\section{Objections and replies}

\begin{objection}[The minimal self]
Every experience has a first-personal character: it is given as mine. This for-me-ness is constant across all experiences, and it is a permanent self \citep{Zahavi2005}.
\end{objection}

\begin{reply}
For-me-ness, if it is a feature of every experience, is present only when there is an experience, and so exists only as long as the activities do. A feature shared by all the stages of a process is not a thing that persists through their cessation. The same holds of the transcendental ego that PHE-20 finds as an irreducible residue of every reduction: what cannot be bracketed in an act of reflection is present in each such act, not before or after them.
\end{reply}

\begin{objection}[Activities need a subject]
Perceiving, remembering and intending are activities of something. If the self is a process, there is nothing that perceives, and the view is incoherent.
\end{objection}

\begin{reply}
The activities do have a subject: the organism (Assumption~\ref{a:org}). The view denies only that the subject is a permanent self. The organism perceives and remembers; the self is its perceiving and remembering, organized around a self-representation.
\end{reply}

\begin{objection}[Sleep]
In dreamless sleep the activities that constitute the self lapse, yet the same self wakes. So the self persists without its activities.
\end{objection}

\begin{reply}
Dreamless sleep is not the cessation of the integrating activities: memory consolidation and the regulation of arousal continue, and they are among the activities that link a subject's states across time. Where every constituent activity does cease, as perhaps under deep anesthesia, the self does not exist during the interval; the process that resumes is the same process if it is caused by the preserved organization, as a performance resumed after an intermission is the same performance. No self persists through the gap.
\end{reply}

\section{Conclusion}

The work assigned to the self is done by the activities that integrate experience, memory and agency; no bearer distinct from the organism is found or needed; and the organism precedes those activities and can outlive them. There is no permanent self. The self is the integrating process, real but impermanent, and it persists only as long as its constituent activities. You are the organism that undergoes it.

\begin{thebibliography}{9}\small
\bibitem[Dennett(1992)]{Dennett1992} Dennett, D.~C. (1992). The self as a center of narrative gravity. In F. Kessel, P. Cole and D. Johnson (eds.), \emph{Self and Consciousness: Multiple Perspectives}. Erlbaum.
\bibitem[Hume(1739)]{Hume1739} Hume, D. (1739). \emph{A Treatise of Human Nature}. Cited by book, part and section.
\bibitem[Metzinger(2003)]{Metzinger2003} Metzinger, T. (2003). \emph{Being No One: The Self-Model Theory of Subjectivity}. MIT Press.
\bibitem[Siderits(2003)]{Siderits2003} Siderits, M. (2003). \emph{Personal Identity and Buddhist Philosophy: Empty Persons}. Ashgate.
\bibitem[Strawson(2009)]{Strawson2009} Strawson, G. (2009). \emph{Selves: An Essay in Revisionary Metaphysics}. Oxford University Press.
\bibitem[Zahavi(2005)]{Zahavi2005} Zahavi, D. (2005). \emph{Subjectivity and Selfhood: Investigating the First-Person Perspective}. MIT Press.
\end{thebibliography}

\end{document}
